This appendix collects the Pascal programs and program fragments that appear in
the text, together with direct Python 3 and Rust translations.  The translations
preserve the same automata, transition tables, start states, final states, and
input conventions as the Pascal originals.

\section{State Transition Function}

\subsection*{Original Pascal}
\begin{verbatim}
type
   Sigma = 'a'..'c';
   State = (s0,s1,s2);
var
   TransitionTable=array[State,Sigma] of State;

function Delta(S:State;A:Sigma):State;
begin
   Delta := TransitionTable [S, A]
end; {Delta}
\end{verbatim}

\subsection*{Python 3}
\begin{verbatim}
from enum import Enum
from collections.abc import Mapping

class State(Enum):
    S0 = 0
    S1 = 1
    S2 = 2

TransitionTable = Mapping[tuple[State, str], State]

def delta(transition_table: TransitionTable,
          state: State,
          letter: str) -> State:
    return transition_table[(state, letter)]
\end{verbatim}

\subsection*{Rust}
\begin{verbatim}
#[derive(Copy, Clone, Debug, Eq, PartialEq)]
enum State {
    S0,
    S1,
    S2,
}

fn state_index(state: State) -> usize {
    match state {
        State::S0 => 0,
        State::S1 => 1,
        State::S2 => 2,
    }
}

fn sigma_index(letter: char) -> usize {
    match letter {
        'a' => 0,
        'b' => 1,
        'c' => 2,
        _ => panic!("letter outside Sigma"),
    }
}

fn delta(transition_table: &[[State; 3]; 3],
         state: State,
         letter: char) -> State {
    transition_table[state_index(state)][sigma_index(letter)]
}
\end{verbatim}

\section{Extended State Transition Function}

\subsection*{Original Pascal}
\begin{verbatim}
const
MaxWordLength = 255; {an arbitrary constraint}
type
   Word = record
      Length :0..MaxWordLength;
      Letters:packed array [0..MaxWordLength] of Sigma
   end; {Word}
function DeltaBar(S:State; W:Word) : State;
{uses the function Delta defined previously}
var
   T:State;
   I:0..MaxWordLength;
begin
   T := S;
   if W.Length>0
      then
         for I := 1 to W.Length do
            T := Delta(T, W.Letters[I]);
   DeltaBar := T
end; {DeltaBar}
\end{verbatim}

\subsection*{Python 3}
\begin{verbatim}
MAX_WORD_LENGTH = 255

def delta_bar(transition_table: TransitionTable,
              state: State,
              word: str) -> State:
    if len(word) > MAX_WORD_LENGTH:
        raise ValueError("word exceeds MaxWordLength")
    current = state
    for letter in word:
        current = delta(transition_table, current, letter)
    return current
\end{verbatim}

\subsection*{Rust}
\begin{verbatim}
const MAX_WORD_LENGTH: usize = 255;

fn delta_bar(transition_table: &[[State; 3]; 3],
             state: State,
             word: &str) -> State {
    if word.chars().count() > MAX_WORD_LENGTH {
        panic!("word exceeds MaxWordLength");
    }
    let mut current = state;
    for letter in word.chars() {
        current = delta(transition_table, current, letter);
    }
    current
}
\end{verbatim}

\section{Acceptance Test}

\subsection*{Original Pascal}
\begin{verbatim}
function Accept(W:Word):Boolean;
{returns TRUE iff W is accepted by the DFA}
begin
   Accept := DeltaBar(s0, W) in FinalState
end; {Accept}
\end{verbatim}

\subsection*{Python 3}
\begin{verbatim}
from collections.abc import Set

def accept(transition_table: TransitionTable,
           final_state: Set[State],
           word: str) -> bool:
    return delta_bar(transition_table, State.S0, word) in final_state
\end{verbatim}

\subsection*{Rust}
\begin{verbatim}
fn accept(transition_table: &[[State; 3]; 3],
          final_state: &[State],
          word: &str) -> bool {
    final_state.contains(&delta_bar(transition_table, State::S0, word))
}
\end{verbatim}

\section{DFA Emulator}

\subsection*{Original Pascal}
\begin{verbatim}
program DFA(input, output);
{This program tests whether input strings are accepted by the     }
{automaton displayed in Figure 1.9. The program expects input from}
{the keyboard, delimited by a carriage return. No error checking  }
{is done; letters outside ['a' .. 'c'] cause a range error.       }
type
    Sigma = 'a'..'c';
    State = (s0, s1, s2);
var
    TransitionTable : array [State, Sigma] of State;
    FinalState      : set of State;
function Delta(s : State; c : Sigma) : State;
begin
    Delta := TransitionTable[s,c]
end; { Delta }
function DeltaBar(s : State) : State;
var
    t : State;
begin
    t := s;
    { Step through the keyboard input one letter at a time. }
    while not eoln(input) do
        begin
            t := Delta(t, input^);
            get(input)
        end;
    DeltaBar := t
end; { DeltaBar }
function Accept : boolean;
begin
    Accept := DeltaBar(s0) in FinalState
end; { Accept }
procedure Initialize;
begin
    FinalState := [s2];
    { Set up the state transition table. }
    TransitionTable [s0,'a'] := s1; TransitionTable [s0,'b'] := s0;
    TransitionTable [s0,'c'] := s2; TransitionTable [s1,'a'] := s2;
    TransitionTable [s1,'b'] := s0; TransitionTable [s1,'c'] := s0;
    TransitionTable [s2,'a'] := s0; TransitionTable [s2,'b'] := s0;
    TransitionTable [s2,'c'] := s1;
end; { Initialize }
begin { DFA }
    Initialize;
    if Accept then
        writeln(output, 'Accepted')
    else
        writeln(output, 'Rejected')
end. { DFA }
\end{verbatim}

\subsection*{Python 3}
\begin{verbatim}
from enum import Enum
import sys

class State(Enum):
    S0 = 0
    S1 = 1
    S2 = 2

TRANSITION_TABLE = {
    (State.S0, "a"): State.S1,
    (State.S0, "b"): State.S0,
    (State.S0, "c"): State.S2,
    (State.S1, "a"): State.S2,
    (State.S1, "b"): State.S0,
    (State.S1, "c"): State.S0,
    (State.S2, "a"): State.S0,
    (State.S2, "b"): State.S0,
    (State.S2, "c"): State.S1,
}
FINAL_STATE = {State.S2}

def delta(state: State, letter: str) -> State:
    return TRANSITION_TABLE[(state, letter)]

def delta_bar(state: State, word: str) -> State:
    current = state
    for letter in word:
        current = delta(current, letter)
    return current

def accept(word: str) -> bool:
    return delta_bar(State.S0, word) in FINAL_STATE

def main() -> None:
    word = sys.stdin.readline().rstrip("\n")
    print("Accepted" if accept(word) else "Rejected")

if __name__ == "__main__":
    main()
\end{verbatim}

\subsection*{Rust}
\begin{verbatim}
use std::io::{self, Read};

#[derive(Copy, Clone, Debug, Eq, PartialEq)]
enum State {
    S0,
    S1,
    S2,
}

fn delta(state: State, letter: char) -> State {
    match (state, letter) {
        (State::S0, 'a') => State::S1,
        (State::S0, 'b') => State::S0,
        (State::S0, 'c') => State::S2,
        (State::S1, 'a') => State::S2,
        (State::S1, 'b') => State::S0,
        (State::S1, 'c') => State::S0,
        (State::S2, 'a') => State::S0,
        (State::S2, 'b') => State::S0,
        (State::S2, 'c') => State::S1,
        _ => panic!("letter outside ['a'..'c']"),
    }
}

fn delta_bar(state: State, word: &str) -> State {
    let mut current = state;
    for letter in word.chars() {
        current = delta(current, letter);
    }
    current
}

fn accept(word: &str) -> bool {
    delta_bar(State::S0, word) == State::S2
}

fn main() {
    let mut input = String::new();
    io::stdin().read_to_string(&mut input).unwrap();
    let word = input.lines().next().unwrap_or("");
    if accept(word) {
        println!("Accepted");
    } else {
        println!("Rejected");
    }
}
\end{verbatim}

\section{Hypothetical Use of HALT}

\subsection*{Original Pascal}
\begin{verbatim}
program CHECK;
{ envisioned usage of HALT }
   function HALT: boolean;
   begin
      { marvelous code goes here }
   end { HALT }

begin { CHECK }
if HALT then
   writeln('The program in file data.p will halt')
else
   writeln('The program in file data.p will not halt')
end { CHECK }.
\end{verbatim}

\subsection*{Python 3}
\begin{verbatim}
def halt() -> bool:
    # Hypothetical oracle from the proof; no such total function exists.
    raise NotImplementedError("marvelous code goes here")

def main() -> None:
    if halt():
        print("The program in file data.p will halt")
    else:
        print("The program in file data.p will not halt")

if __name__ == "__main__":
    main()
\end{verbatim}

\subsection*{Rust}
\begin{verbatim}
fn halt() -> bool {
    // Hypothetical oracle from the proof; no such total function exists.
    unimplemented!("marvelous code goes here")
}

fn main() {
    if halt() {
        println!("The program in file data.p will halt");
    } else {
        println!("The program in file data.p will not halt");
    }
}
\end{verbatim}

\section{Contradictory Test Program}

\subsection*{Original Pascal}
\begin{verbatim}
program TEST;
{ to be placed in the file data.p }
var FOREVER: boolean;
   function HALT: boolean;
   begin
      { marvelous code goes here }
   end; { HALT }

begin { TEST }
FOREVER := false;
if HALT then
   repeat
      FOREVER := false;
   until FOREVER
else
   writeln('This program halts')
end { TEST }.
\end{verbatim}

\subsection*{Python 3}
\begin{verbatim}
def halt() -> bool:
    # Hypothetical oracle from the proof; no such total function exists.
    raise NotImplementedError("marvelous code goes here")

def main() -> None:
    forever = False
    if halt():
        while not forever:
            forever = False
    else:
        print("This program halts")

if __name__ == "__main__":
    main()
\end{verbatim}

\subsection*{Rust}
\begin{verbatim}
fn halt() -> bool {
    // Hypothetical oracle from the proof; no such total function exists.
    unimplemented!("marvelous code goes here")
}

fn main() {
    let mut forever = false;
    if halt() {
        while !forever {
            forever = false;
        }
    } else {
        println!("This program halts");
    }
}
\end{verbatim}
